\documentclass[journal]{IEEEtran}

\usepackage{amsmath,amssymb}
\usepackage{booktabs}
\usepackage{graphicx}
\usepackage{hyperref}
\usepackage{xcolor}
\usepackage{placeins}
\usepackage{balance}
\hypersetup{hidelinks}

\definecolor{draftred}{RGB}{150,20,20}
\newcommand{\draftwarning}{\textcolor{draftred}{DRAFT---NOT FOR SUBMISSION}}
\newcommand{\MVTecSpatialAUROC}{TBD}
\newcommand{\MVTecPixelAUROC}{TBD}
\newcommand{\WireNNAUROC}{TBD}
\newcommand{\WireNNRecall}{TBD}
\newcommand{\WorstViewRecall}{TBD}
\newcommand{\StressTopoBFour}{TBD}
\newcommand{\DeployTopoBFour}{TBD}
\newcommand{\DeployGeomBFour}{TBD}
\newcommand{\DeployDeltaBFour}{TBD}
\newcommand{\DeployDeltaBSix}{TBD}
\newcommand{\DeployDeltaBEight}{TBD}
\newcommand{\DeployDeltaBTen}{TBD}
\newcommand{\DeployBFourHolm}{TBD}
\newcommand{\MuJoCoRuns}{TBD}
\newcommand{\MuJoCoFOVmm}{TBD}
\newcommand{\MuJoCoIKmm}{TBD}
\newcommand{\AuditChecks}{TBD}
\IfFileExists{../../artifacts/papers/inspectbot/final/paper_assets/inspect_numbers.tex}
  {\input{../../artifacts/papers/inspectbot/final/paper_assets/inspect_numbers.tex}}{}

\title{Topology-Risk Active Inspection for Wire Harnesses:\\
Public Cable Defects, Paired Budget Tests, and Executable Robot Scans}
\author{Anonymous Authors}
\date{}

\begin{document}
\maketitle

\begin{center}
\draftwarning
\end{center}

\begin{abstract}
Robot inspection of branched wire harnesses couples anomaly detection with the
choice of where to spend a limited scan budget. We present InspectBot, a hybrid
benchmark grounded in two public real-image datasets and executable 3-D robot
scans. A spatial Gaussian model trained only on normal MVTec cable images reaches
\MVTecSpatialAUROC\% image and \MVTecPixelAUROC\% pixel AUROC. On independent
stripped-wire data, a spatial nearest-neighbor reference reaches
\WireNNAUROC\% AUROC but only \WireNNRecall\% recall at a normal-only threshold.
Occlusion plus darkening reduces anomaly recall to \WorstViewRecall\%, exposing
a perception floor that no scan policy repairs. Under a disjoint 2,000-scenario
protocol, fixed topology-risk ordering improves criticality-weighted recall over
geometry coverage by \DeployDeltaBFour, \DeployDeltaBSix, \DeployDeltaBEight,
and \DeployDeltaBTen{} percentage points at budgets 4, 6, 8, and 10; all remain
significant after a 30-test Holm correction. The advantage disappears at full
coverage and reverses under naive revisits. Direct MuJoCo execution completes
all \MuJoCoRuns{} runs with sub-\MuJoCoFOVmm{}-mm field-of-view and
sub-\MuJoCoIKmm{}-mm visual-IK error. The evidence supports risk-aware budget
allocation, not learned topology, hardware safety, or factory performance.
\end{abstract}

\noindent\textit{Note to Practitioners---}
Wire-harness inspection stations often have less camera time than is needed to
view every terminal, branch, label, and clip at equal quality. This paper tests
a simple operational principle: spend the limited scan budget first at sites
whose process role makes a missed defect costly, but stop trusting the planner
when image quality is too poor. On public cable and stripped-wire images, the
experiments show that adverse views can impose a perception floor that no view
ordering policy repairs. Under cleaner views, topology-risk ordering improves
criticality-weighted recall at partial coverage, but the advantage disappears
at full coverage and reverses under poorly designed revisits. The current robot
is a digital twin, so the reported rates are not factory escape rates. Before
deployment, practitioners must calibrate the actual cameras and fixtures,
measure cycle time, and validate the policy across suppliers, lighting and real
automotive harness defects.

\begin{IEEEkeywords}
Active perception, anomaly detection, industrial inspection, robot vision,
wire-harness assembly.
\end{IEEEkeywords}

\section{Introduction}

Wire-harness inspection is not uniform search.  Damage near a branch junction,
high-current path, or safety-critical connector can carry greater engineering
risk than an equally visible defect on a redundant segment.  At the same time,
the camera and anomaly detector may be least reliable at occluded or reflective
locations.  An inspection robot must therefore decide not merely whether an
image is anomalous, but which site should be observed next and whether revisiting
it is preferable to covering an unobserved site.

Industrial anomaly detection has strong controlled benchmarks and capable
one-class methods \cite{bergmann2019mvtec,defard2021padim,roth2022patchcore}.
CableInspect-AD shows that expert cable anomalies remain difficult outside
controlled factories \cite{arodi2024cableinspect}.  Wire-harness reviews
similarly identify robustness, practical validation, and exploitation of
intrinsic harness structure as open problems \cite{wang2024review}.  Existing
evidence, however, rarely couples a real-image detector to paired resource-
allocation tests and executable scan motions.  Merely rendering a robot beside
a cable does not close that gap.

InspectBot narrows the claim to a falsifiable question: given frozen detector
scores and a fixed engineering-risk map, does topology-aware allocation recover
more risk-weighted anomaly mass than non-topological policies at the same number
of observations?  Its contributions are:

\begin{enumerate}
  \item two integrity-audited public-data experiments with normal-only fitting
  and threshold calibration, including an independent stripped-wire refit;
  \item a 600-row real-image score cache under four deterministic appearances,
  exposing both a useful clean-view regime and a severe adverse-view failure;
  \item paired stress and deployment protocols covering 126,000 active-policy
  rows, exact frozen seeds, bootstrap intervals, multiplicity-corrected tests,
  an explicit shuffled-risk control, and a corrected offline oracle;
  \item a direct MuJoCo execution audit that preserves all allocation outcomes
  while testing bounded 3-D motion to the 12 inspection sites; and
  \item a machine-audited claim boundary separating real images, synthetic
  appearance transforms, simulated robot motion, and absent hardware evidence.
\end{enumerate}

\begin{figure*}[t]
  \centering
  \includegraphics[width=0.96\textwidth]{%
    ../../artifacts/papers/inspectbot/final/paper_assets/system_overview.pdf}
  \caption{InspectBot separates allocation evidence from execution evidence.
  Left: the UR5e visual twin executes bounded scanner motions in MuJoCo.  Right:
  12 frozen inspection sites and their engineering-risk weights; the displaced
  dotted node is a process checkpoint physically colocated with the junction.
  Detector scores come from held-out real images, not from the rendered camera.}
  \label{fig:overview}
\end{figure*}

\section{Related Work}

\paragraph{Industrial anomaly detection.}
MVTec AD established a controlled multi-category one-class benchmark
\cite{bergmann2019mvtec}.  PaDiM fits a Gaussian distribution at each pretrained
feature location \cite{defard2021padim}; PatchCore instead retains a coreset of
nominal patch features \cite{roth2022patchcore}.  We use these as transparent
method families, not as a novelty claim.  Our compact spatial Gaussian keeps 64
deterministic channels from a frozen ResNet-18, while the independent validation
also includes a position-wise nearest-neighbor reference.

CableInspect-AD contains expert-annotated high-resolution power-line cable images
with multiple real-world anomaly types \cite{arodi2024cableinspect}.  We audited
its official code but could not obtain the dataset from its official endpoints
without an access request; consequently no CableInspect number appears in our
tables.  The FAU/FAPS Stripped Wire Dataset contains good, cut-strand, and
pulled-strand images \cite{scheck2025stripped}.  Its associated PB-IAD study uses
semantic foundation-model prompts and a PatchCore comparison
\cite{hofmann2025pbiad}.  That comparison selects a threshold using labeled test
data, whereas every InspectBot threshold is selected from normal calibration
images only; thresholded F1 values are therefore not directly comparable.

\paragraph{Wire-harness perception and active inspection.}
Multi-branch 3-D vision has been validated in a real dual-robot assembly setting
\cite{nguyen2024branches}, which is stronger hardware evidence than ours.  We
address a different decision: distributing observations after the sites and a
risk map are known.  This follows the active-perception view that sensing actions
should be selected jointly with task goals \cite{bajcsy1988active}.  The risk map
here is engineering input, not learned from labels.  A shuffled-risk policy is
therefore essential: it tests whether improvements require the correct
site--risk binding rather than merely a nonuniform scan order.

\section{Public-Data Perception}

\subsection{Data integrity and splits}

The MVTec cable audit fixes the transport mirror at commit
\texttt{c75b39616f84} and the extracted tree at SHA-256 prefix
\texttt{5ecf3b6b9fee}.  It contains 224 normal training images, 58 normal test
images, 92 anomalous test images across eight defect types, and 92 masks.
A seed-20260827 permutation assigns 180 normals to fitting and 44 to threshold
calibration.  The public test set is untouched until evaluation.

The Stripped Wire archive is fixed to DOI
\href{https://doi.org/10.5281/zenodo.16686806}{10.5281/zenodo.16686806}, version
v2, MD5 \texttt{9e1944f1b7c9bc9e7db433f7bdcf2694}.  Its 200 released
PatchCore normals are split 160/40 for fitting/calibration; all 133 good, 68
cut-strand, and 99 pulled-strand test images remain held out.  Exact-file hashing
finds no train--test duplicate.  Three optional VLM references are excluded.
Aspect ratio is preserved and padded to $224\!\times\!224$, avoiding a
test-label-dependent crop.

\subsection{Detector and baselines}

Frozen ImageNet ResNet-18 features from layers 1--3 are aligned to a
$28\!\times\!28$ grid.  A seed-fixed subset of 64 channels forms feature
$x_{ijp}$ for normal image $j$ at spatial position $p$.  The spatial Gaussian
score is
\begin{equation}
 a_p(x)=\sqrt{(x_p-\mu_p)^\top(\Sigma_p+10^{-2}I)^{-1}(x_p-\mu_p)},
\end{equation}
and the image score averages the top 1\% of spatial values.  The strong spatial
nearest-neighbor reference replaces the Gaussian distance with the closest
normal feature at the same position.  RGB statistics and globally pooled
ResNet features, each scored with Ledoit--Wolf Mahalanobis distance, are weaker
controls.  Each method's threshold is its own 99th percentile on normal
calibration scores.  No anomalous image selects features, preprocessing,
thresholds, or models.

\begin{table}[t]
  \centering
  \caption{MVTec cable test results (58 normal/92 anomaly).  Values are percent;
  the decision threshold uses 44 normal calibration images only.}
  \label{tab:mvtec}
  \resizebox{\columnwidth}{!}{%
    \input{../../artifacts/papers/inspectbot/final/paper_assets/mvtec_perception_table.tex}}
\end{table}

\begin{table}[t]
  \centering
  \caption{Independent Stripped Wire refit (133 normal/167 anomaly).  Spatial
  NN is an exploratory strong reference, not the prespecified candidate.}
  \label{tab:wire}
  \resizebox{\columnwidth}{!}{%
    \input{../../artifacts/papers/inspectbot/final/paper_assets/wire_perception_table.tex}}
\end{table}

\begin{figure*}[t]
  \centering
  \includegraphics[width=0.97\textwidth]{%
    ../../artifacts/papers/inspectbot/final/paper_assets/perception_roc.pdf}
  \caption{Image-level ROC curves recomputed from frozen per-image CSV files.
  MVTec favors the spatial Gaussian model; the independent wire-end task favors
  the spatial nearest-neighbor reference.}
  \label{fig:roc}
\end{figure*}

\subsection{Perception results}

On MVTec cable, the spatial Gaussian reaches \MVTecSpatialAUROC\% image AUROC
(stratified image-bootstrap 95\% CI 87.3--96.1) and 94.7\% AUPR.  Its
normal-calibrated operating point gives 65.2\% recall, 95.2\% precision, and
5.2\% false-positive rate.  Pixel AUROC is \MVTecPixelAUROC\%, but pixel AUPR is
48.2\% and F1 is 52.5\%; localization is useful but far from solved.

On Stripped Wire, spatial nearest neighbor is best at \WireNNAUROC\% AUROC
(76.5--86.1) and 83.9\% AUPR.  Spatial Gaussian reaches 78.3\% AUROC, exceeding
RGB statistics by 10.3 points (paired bootstrap CI 4.7--16.0) and global
features by 4.8 points (0.2--9.4), but it is 3.2 points below spatial NN
($-5.4$ to $-1.1$).  The strict spatial-NN threshold yields only
\WireNNRecall\% recall.  Its per-defect AUROC is 71.1\% on cut strands and
88.6\% on pulled strands, exposing the harder failure mode.  We therefore do
not claim method superiority or production readiness.

\begin{figure*}[t]
  \centering
  \includegraphics[width=0.97\textwidth]{%
    ../../artifacts/papers/inspectbot/final/paper_assets/qualitative_anomalies.pdf}
  \caption{Deterministic qualitative audit.  The three largest MVTec defect
  classes by test count and both Stripped Wire defect classes are shown; within
  each class the image nearest its median anomaly score is selected.  Per-image
  robust color scaling aids localization display but is not used for decisions.
  Boundary responses on wire ends reveal a concrete localization weakness.}
  \label{fig:qualitative}
\end{figure*}

\section{Hybrid Active Inspection}

\subsection{Real-image score bridge}

Every held-out MVTec image is scored under four deterministic appearances:
clean; dark plus blur; rotation plus contrast; and deterministic occlusion plus
darkening.  Each appearance $v$ is independently calibrated on the same 44
normal images.  We store
\begin{equation}
 q_{iv}=\frac{s_{iv}-\tau_v}{1.4826\,\mathrm{MAD}_v},
\end{equation}
where $\tau_v$ is the 99th-percentile normal threshold.  A site's prediction is
the sign of the mean of all observed $q_{iv}$.  These transforms are controlled
stress tests, not physical views or camera trajectories.

\begin{table}[t]
  \centering
  \caption{Frozen detector robustness by appearance (150 held-out images).
  Each threshold is calibrated on transformed normal training images.}
  \label{tab:views}
  \resizebox{\columnwidth}{!}{%
    \input{../../artifacts/papers/inspectbot/final/paper_assets/view_robustness_table.tex}}
\end{table}

The clean and dark-blur appearances retain 92.2\% and 92.6\% AUROC.  Rotation
plus contrast falls to 62.9\%; occlusion plus darkening falls to 55.2\% AUROC
and \WorstViewRecall\% recall.  This observation motivates two frozen protocols
rather than deletion of the failed setting.

\subsection{Sites, outcome, and policies}

Each paired scenario assigns a held-out real image to each of 12 sites on a
branched topology, with anomaly probability 0.30 conditioned on at least one
normal and one anomalous site.  Defect type is sampled uniformly before an
image is selected within type.  Site risk $r_i$ and defect-severity coefficient
$d_i$ are fixed before the run.  For true anomaly indicator $y_i$ and final
prediction $\hat y_i$, the primary metric is
\begin{equation}
 \mathrm{CWR}=\frac{\sum_i r_i d_i y_i\hat y_i}
                    {\sum_i r_i d_i y_i}.
\end{equation}
This metric intentionally values critical sites.  Ordinary recall, precision,
F1, false-positive rate, coverage, revisits, travel, and complete critical
success are secondary outcomes.

The seven policies are raster order, random scanning, nearest-unvisited geometry
coverage, coverage-first uncertainty revisit, topology risk plus uncertainty,
a shuffled-risk negative control, and an offline label/score-aware oracle.  The
topology policy selects unvisited site $i$ approximately by
$r_i/(1+\mathrm{travel}_i)$ and, after coverage, combines risk with score
uncertainty.  It is a deterministic engineering heuristic, not a learned
planner.  The oracle solves an exact multiple-choice knapsack over per-site view
prefixes and is descriptive only.

\subsection{Frozen paired protocols and statistics}

Stress-v1 uses 1,000 scenarios (seeds 77,000,000--77,000,999), budgets
$\{4,6,8,10,12,16\}$, and the appearance order adverse-dark, dark-blur,
rotate-contrast, clean.  Deployment-v2 was written after retaining that result;
it changes only to clean-first order and uses 2,000 disjoint scenarios (seeds
78,000,000--78,001,999).  Detector scores, topology, risks, severities, policy
equations, thresholds, budgets, and tests remain fixed.

At each budget we compare topology risk with five non-oracle baselines using
paired, one-sided Wilcoxon signed-rank tests on CWR.  All 30 $p$-values are Holm
adjusted together.  Scenario-paired bootstrap intervals use 10,000 resamples.
Budgets 4--10 are confirmatory under-coverage regimes; budget 12 is a
prespecified coverage equality check and budget 16 is exploratory.  An initial
greedy implementation mislabeled as ``oracle'' was not a true upper bound.  We
retain that run, correct only the descriptive oracle, and audit that all 72,000
non-oracle rows and the complete paired-comparison CSV remain unchanged.

\begin{figure*}[t]
  \centering
  \includegraphics[width=0.97\textwidth]{%
    ../../artifacts/papers/inspectbot/final/paper_assets/active_budget.pdf}
  \caption{Criticality-weighted recall versus scan budget on shared 0--75\%
  axes.  Stress-v1 is a retained negative result: adverse-first observations
  impose a perception floor.  Clean-first deployment-v2 reveals a topology-risk
  allocation advantage under partial coverage, followed by equality and a
  late-revisit reversal.  Bands are scenario-bootstrap 95\% intervals.}
  \label{fig:budget}
\end{figure*}

\begin{table*}[t]
  \centering
  \caption{Deployment-v2 criticality-weighted recall in percent.  Differences
  are paired topology minus geometry.  Holm $p$ values belong to the conservative
  family of all 30 topology-versus-baseline comparisons; budget 16 used a
  one-sided superiority alternative and is exploratory.}
  \label{tab:active}
  \resizebox{0.82\textwidth}{!}{%
    \input{../../artifacts/papers/inspectbot/final/paper_assets/active_budget_table.tex}}
\end{table*}

\section{Active-Inspection Results}

\paragraph{The adverse-first protocol fails informatively.}
Topology risk reaches only \StressTopoBFour\% CWR at budget 4, and no
non-oracle policy shows a stable advantage through budget 12: the first site
observation has only \WorstViewRecall\% anomaly recall.  Random revisits
occasionally encounter a less damaging transform earlier.  Allocation cannot
compensate for an arbitrarily poor sensing regime.

\paragraph{Topology matters under clean partial coverage.}
In deployment-v2, topology risk reaches \DeployTopoBFour\% CWR at budget 4,
versus \DeployGeomBFour\% for geometry, a \DeployDeltaBFour-point paired gain
(95\% CI 8.0--10.9, Holm-adjusted $p$ value \DeployBFourHolm).  The gains at budgets 6, 8, and
10 are \DeployDeltaBSix, \DeployDeltaBEight, and \DeployDeltaBTen{} points.  At
all four confirmatory budgets topology also exceeds raster, random, uncertainty,
and shuffled-risk after the same 30-test correction.  The shuffled control
supports the interpretation that correct site--risk binding, rather than an
arbitrary nonuniform order, drives the primary metric.

\paragraph{Coverage and revisits bound the claim.}
At budget 12, all coverage-first policies inspect every site once and tie at
64.8\% CWR.  At budget 16, topology risk falls to 61.2\%, below geometry's
63.6\% and uncertainty's 62.6\%.  Averaging later degraded-view scores can turn
a previously detected anomaly negative; risk-weighted revisiting then harms the
same high-value sites it tries to protect.  The corrected offline oracle uses an
average of 2.64 scans and reaches 66.5\% CWR at budget 4, 68.2\% at 6, and
68.4\% from budget 8 onward.  The gap suggests a future learned stopping rule,
but oracle knowledge cannot be deployed.

\section{Direct MuJoCo Audit}

The direct audit uses seeds 78,000,100--78,000,119, four executable policies,
and budgets 4, 8, 12, and 16.  At each decision, the current 3-D inspection-site
pose is read from MuJoCo \cite{todorov2012mujoco}.  The scanner moves in
increments no larger than 25 mm, with four physics steps per increment, and a
scan is accepted only within 3 mm of the site.  A detailed UR5e plus dual-lens
head follows the task-space scanner through inverse kinematics.  The arm mesh is
collision-disabled so it cannot silently perturb the cable benchmark.

\begin{table}[t]
  \centering
  \caption{Direct execution diagnostics pooled over four budgets per policy.
  Travel is world-space scanner-path length.}
  \label{tab:mujoco}
  \resizebox{\columnwidth}{!}{%
    \input{../../artifacts/papers/inspectbot/final/paper_assets/mujoco_execution_table.tex}}
\end{table}

All \MuJoCoRuns{} runs reach every prescribed site.  Maximum accepted
field-of-view error is \MuJoCoFOVmm{} mm and maximum visual IK error is
\MuJoCoIKmm{} mm; every simulation state remains finite.  For identical seed,
policy, and budget, 14 shared perception/allocation fields exactly match
deployment-v2.  This is deliberate: the new evidence is executable 3-D motion,
not a second performance sample.  The detector scores are replayed from public
images rather than computed from MuJoCo renders, and neither collision forces
nor real camera calibration is validated.

\section{Limitations}

The primary metric multiplies an engineered risk map by defect severity, and
the topology policy receives that same risk map.  Paired geometry, uncertainty,
random, raster, and shuffled-risk controls establish resource-allocation value,
but they do not show that risk was learned or optimally chosen.  An actual
factory must derive weights through traceable hazard analysis and test their
sensitivity.

The active benchmark repeatedly samples a finite 150-image MVTec test pool.
Intervals quantify scenario resampling, not unseen cables, suppliers, cameras,
or factories.  Deterministic transforms are appearance perturbations rather
than robot viewpoints; site localization, autofocus, exposure control, depth,
and temporal registration are absent.  Stripped Wire is an independent refit,
not zero-shot domain transfer, and both public datasets show controlled cable
cross-sections or single wire ends rather than a complete automotive harness.

Only one frozen backbone/channel seed is evaluated.  The study does not claim a
state-of-the-art anomaly detector, an optimal active policy, or superiority over
CableInspect-AD and PB-IAD.  Deployment-v2 was designed after observing the
stress-v1 perception floor, although its new protocol and disjoint seeds were
frozen before execution.  Reporting both versions is mandatory.

Finally, MuJoCo validates bounded task-space reachability with a visual IK twin.
It does not validate joint control, collision avoidance, cable--robot contact,
force limits, camera measurements, timing, safety interlocks, maintenance, or
production economics.  Hardware experiments and a prospective cross-factory
protocol are required before any deployment claim.

\section{Ethics and Safety}

The used datasets contain industrial objects and no personal data.  MVTec cable
is CC BY-NC-SA 4.0; Stripped Wire v2 is CC BY 4.0.  A release should distribute
download instructions and hashes rather than silently relicense source images.
Automated inspection may reduce repetitive visual work and improve traceability,
but it may also displace tasks or create automation bias.  A missed electrical
defect can be safety-critical.  The present detector and simulator must not be
used to waive human inspection or safety validation.

\section{Reproducibility}

The research package includes source URLs, transport revision, archive and tree
hashes, fit/calibration manifests, checkpoints, all per-image predictions and
maps, the 600-row view cache, per-scenario active CSVs, the retained oracle bug,
frozen protocols, MJCF, per-run motion traces, plot code, and a one-command paper
build.  An independent audit recomputes public counts, file hashes, AUROC/AUPR,
threshold outcomes, every active aggregate, oracle upper-bound behavior, all
non-oracle correction rows, MuJoCo tolerances, and cross-system outcome identity.
All \AuditChecks{} checks must pass before the release PDF is copied.

\balance
\bibliographystyle{plain}
\bibliography{references}

\end{document}
